% Generated draft: requires mathematical and source review.
\documentclass[11pt]{article}
\usepackage[margin=25mm]{geometry}
\usepackage{fontspec}
\setmainfont{Latin Modern Roman}
\usepackage{amsmath,amssymb,mathtools,unicode-math}
\setmathfont{Latin Modern Math}
\usepackage{enumitem,needspace,xurl,hyperref}
\setlength{\parindent}{0pt}
\setlength{\parskip}{5pt}
\setlength{\emergencystretch}{3em}
\newcommand{\N}{\mathbb{N}}
\newcommand{\Z}{\mathbb{Z}}
\newcommand{\Q}{\mathbb{Q}}
\newcommand{\R}{\mathbb{R}}
\newcommand{\C}{\mathbb{C}}
\newcommand{\F}{\mathbb{F}}
\newcommand{\A}{\mathbb{A}}
\newcommand{\PP}{\mathbb{P}}
\newcommand{\E}{\mathbb{E}}
\newcommand{\eps}{\varepsilon}
\newcommand{\dd}{\,\mathrm{d}}
\begin{document}
\section*{UnsolvedMath ID: 801147}
\textbf{Source problem code:} catalogue entry
\subsection{Definitions and mathematical statement}
A separable \(\mathrm{II}_1\) factor is a von Neumann factor with separable predual and a faithful normal tracial state \(\tau\) normalized by \(\tau(1)=1\). Write \(\|x\|_{2,\tau}=\tau(x^*x)^{1/2}\). Given an ultrafilter \(\mathcal U\) on an index set \(K\), its tracial ultrapower is
\[
N^{\mathcal U}=\ell^\infty(K,N)\big/\{(x_k)_{k\in K}:\lim_{k\to\mathcal U}\|x_k\|_{2,\tau}=0\}.
\]
Here \(\ell^\infty(K,N)\) consists of operator-norm bounded families. The source uses a common \(\mathcal U\) but does not specify its index set, nonprincipality, or quantifier; these remain unresolved parameters of the question.

Interpret \(\ast\) as the reduced tracial free product over the scalar algebra, with respect to the normalized traces. For tracial von Neumann algebras \(A,B\), write \(A\equiv B\) when every sentence in the usual continuous first-order language of tracial von Neumann algebras has the same value in \(A\) and \(B\).

\textbf{Question (Elayavalli, Problem 49).} Suppose \(I,J\) are finite sets and \((N_i)_{i\in I}\), \((M_j)_{j\in J}\) are families of separable \(\mathrm{II}_1\) factors. Does
\[
\underset{i\in I}{\ast}\bigl(N_i^{\mathcal U}\bigr)
\cong
\underset{j\in J}{\ast}\bigl(M_j^{\mathcal U}\bigr)
\]
imply that there is a bijection \(\sigma:I\to J\) such that
\[
N_i\equiv M_{\sigma(i)}\qquad\text{for every }i\in I?
\]
The displayed isomorphism is understood as an isomorphism of tracial von Neumann algebras. The ultrapowers are taken separately before forming the free products. The conclusion concerns elementary equivalence of the original factors.

\textbf{Unresolved qualifications.} The source says only that the index sets are finite, so empty sets are not expressly excluded; it supplies no convention for an empty free product. Thus the empty-index cases require clarification. The interpretation of free product and elementary equivalence above is the usual tracial interpretation, but the supplied excerpt does not explicitly define these conventions. No choice of ultrafilter scope or additional nonemptiness hypothesis is asserted here as part of the source.
\subsection{Short English statement}
Elayavalli asks whether an isomorphism between finite free products of tracial ultrapowers determines the number of free factors and matches the original factors up to elementary equivalence. The ultrafilter assumptions and empty-index cases need clarification. The registry labels the problem open; this is a supplied status, not an independently verified assessment.
\subsection{Sources}
\begin{itemize}
\item Supplied local UnsolvedMath registry record, permanent ID \(801147\), exact source problem code \texttt{catalogue entry}, title ``Ultraproduct \(\mathrm{II}_1\) Factors --- Problem 49.'' The record attributes the problem to Elayavalli, with proposed year 2025, and labels it open.
\item Supplied source attribution: Elayavalli, \emph{50 Open Problems: Ultraproduct \(\mathrm{II}_1\) Factors} (2025), Problem 49. Supplied source URL: \url{https://arxiv.org/abs/2511.20377}. The background records access on 2026-10-16 and extraction by \texttt{arxiv\_latex}; these are supplied provenance details, not a claim of independent access or verification.
\item Supplied upstream revision: \texttt{d371ee917ac9647fbb06cd6fc529deec6f3dda82}. The registry record was created and updated on \texttt{2026-10-16T14:00:00Z}.
\end{itemize}
\end{document}
